\documentclass[11pt,a4paper,oneside]{book}
\usepackage{mathscommun}
\setniveau{Première}
\renewcommand{\ReferenceProgramme}{Spécialité mathématiques -- programme applicable en 2026--2027}
\renewcommand{\EditionDocument}{Édition du 23 octobre 2026}
\hypersetup{pdftitle={Première spécialité - Guide et couverture}}
\begin{document}
\titrecollection{Guide, couverture et révisions}{Organiser le travail et préparer le supérieur.}{Programme officiel, progression, dépendances des preuves, automatismes corrigés et parcours vers un master de mathématiques.}
\frontmatter\tableofcontents\mainmatter
\chapter{Périmètre, sources et progression}
\section{Le programme de référence}
Ce volume couvre la \textbf{spécialité mathématiques de première générale},
programme publié au \textbf{BO n\textsuperscript{o} 14 du 2 avril 2026}, arrêté
du 26 février 2026, NOR \textbf{MENE2602917A}, applicable à la rentrée 2026--2027.
La référence est le texte et son annexe, consultés le 23 octobre 2026 :
\href{https://www.education.gouv.fr/bo/2026/Hebdo14/MENE2602917A}{texte officiel et annexe du programme}.
Le portail
\href{https://eduscol.education.gouv.fr/5817/programmes-et-ressources-en-mathematiques-voie-gt}{Éduscol : programmes et ressources en mathématiques}
permet de retrouver les programmes de la voie générale et technologique.

Le périmètre ne désigne pas l'enseignement commun des élèves n'ayant pas
choisi la spécialité, ni les programmes des séries technologiques.
L'objectif personnel de master motive des compléments explicites ; il ne
modifie pas le texte officiel et ne rend pas ces compléments obligatoires au lycée.
Les formulations ci-dessous sont des descriptions pédagogiques synthétiques,
pas une reproduction intégrale du Bulletin officiel.

\section{Les quatre documents et leurs usages}
Le cours contient 30 chapitres. Le recueil propose 180 exercices, six par
chapitre, et huit problèmes de synthèse ; tous possèdent un corrigé correspondant.
Ce guide ajoute 80 automatismes avec réponses. Les programmes Python sont
exécutables sans bibliothèque externe ; leurs 18 vérifications contrôlent
les calculs et les contrats d'entrée.

Les chapitres 1 à 14 et 16 à 24 portent les thèmes officiels, avec des
approfondissements signalés. Le chapitre 15 et les chapitres 25 à 30 constituent
des compléments. Une notion reprise de seconde est soit rappelée ici, soit
renvoyée au volume de seconde pour sa démonstration initiale.

\needspace{25\baselineskip}
\section{Progression possible}
Un parcours de travail alterne les domaines plutôt que de repousser toutes
les probabilités à la fin.
\begin{longtable}{>{\raggedright\arraybackslash}p{.21\linewidth}>{\raggedright\arraybackslash}p{.69\linewidth}}
\toprule Période indicative & Travail principal et retours\\\midrule
Début & Diagnostic de seconde ; logique et calcul (1--3) ; définition des suites (4--5) ; conditionnement (19).\\
Premier approfondissement & Radian (14), produit scalaire (16), nombre dérivé (8), modèles discrets (6) ; automatismes réguliers.\\
Milieu & Calcul des dérivées (9), variations (10), géométrie métrique (17--18), épreuves répétées (20).\\
Consolidation & Optimisation (11), exponentielle (12--13), limites intuitives des suites (7), variables aléatoires (21--22).\\
Fin du socle & Échantillonnage (23), projets Python (24), huit problèmes de synthèse, reprises ciblées.\\
Enrichissement & Trigonométrie fonctionnelle (15), fondements d'analyse (25--26), géométrie vectorielle et probabiliste (27--28), convergence (29), liens vers le master (30).\\\bottomrule
\end{longtable}
Les exercices de preuve peuvent être répartis sur l'année. Il n'est pas
nécessaire de maîtriser les chapitres 25 à 30 pour commencer les exercices du socle.
Pour chaque séance autonome : une courte réactivation, une notion nouvelle,
un exemple refait sans aide, des exercices, puis une explication orale.

\chapter{Table de couverture du programme}
Les numéros désignent les chapitres du cours. Les exercices portent les mêmes
numéros ; S1 à S8 sont les problèmes de synthèse.
\section{Algèbre et suites}
\begin{longtable}{>{\raggedright\arraybackslash}p{.70\linewidth}>{\raggedright\arraybackslash}p{.20\linewidth}}
\toprule Contenus et capacités & Localisation\\\midrule
Suite comme fonction sur les entiers ; $u_n$, $u(n)$ ; définition explicite, récurrence et algorithme ; motifs géométriques et dénombrement & 4, 24\\
Calcul de termes, représentation par points, construction graphique d'une récurrence, étude des variations & 4, 7\\
Suites arithmétiques et géométriques ; termes généraux ; sommes consécutives ; croissance absolue et proportionnelle & 5, 6 ; S2\\
Sommes $1+\cdots+n$ et $1+q+\cdots+q^n$ ; cas $q=1$ & 5\\
Limite intuitive : finie, infinie, absence de limite ; limites formelles distinguées comme compléments & 7\\
Second degré : formes factorisée, développée, canonique ; racines, signe, somme et produit ; choix de forme & 2, 3\\
Discriminant et formule des racines ; équations et inéquations ; polynômes s'annulant en deux points donnés & 3 ; ex. 3.1--3.5\\
Algorithmes : termes, sommes, seuils, factorielle, Fibonacci, Syracuse avec borne explicite & 4--7, 24 ; module Python\\
Approfondissements proposés : sommes de carrés et cubes, Hanoï, remboursements, facteurs de $X^n-a^n$, cubique avec racine connue & 2, 5, 6\\\bottomrule
\end{longtable}
\section{Analyse}
\begin{longtable}{>{\raggedright\arraybackslash}p{.70\linewidth}>{\raggedright\arraybackslash}p{.20\linewidth}}
\toprule Contenus et capacités & Localisation\\\midrule
Taux de variation, sécantes, nombre dérivé comme limite ; unités, vitesse instantanée, coût marginal & 8\\
Tangente : pente, lecture et construction, équation ; approximation affine et contrôle d'erreur & 8, 11 ; guide\\
Dérivabilité sur un intervalle ; dérivées carré, cube, inverse, racine ; puissances entières & 8, 9\\
Dérivées de somme, produit, inverse et quotient ; expressions composées avec une fonction affine & 9, 12\\
Non-dérivabilité de la valeur absolue en zéro ; racine carrée sans dérivée finie en zéro & 8\\
Parité et symétries des graphes ; signe de dérivée, monotonie, constance, extrémums ; lire une représentation & 10 ; guide\\
Optimisation géométrique et numérique, inégalités, comparaison de courbes & 3, 10, 11 ; S1\\
Exponentielle définie par $f'=f$, $f(0)=1$ ; $e^x$, loi d'addition, valeurs opposées, positivité, croissance et courbe & 12\\
Fonctions $e^{at}$, modèles de croissance et décroissance ; lien avec suites géométriques & 12--13 ; S2\\
Algorithmes : liste de taux, Newton favorable, Euler, approximation de $e$ par $(1+1/n)^n$ & 8, 11, 13, 24, 26, 29\\
Cercle trigonométrique, enroulement, radian et longueur d'arc ; sinus et cosinus réels, lien avec le triangle rectangle & 14\\
Valeurs remarquables et angles associés ; algorithme d'Archimède pour $\pi$ & 14, 24\\
Compléments : identités d'addition, dérivées trigonométriques, séries, logarithme & 15, 25--26\\\bottomrule
\end{longtable}
\section{Géométrie}
\begin{longtable}{>{\raggedright\arraybackslash}p{.70\linewidth}>{\raggedright\arraybackslash}p{.20\linewidth}}
\toprule Contenus et capacités & Localisation\\\midrule
Produit scalaire défini par projection signée, expression avec cosinus, orthogonalité & 16\\
Symétrie, bilinéarité, norme, coordonnées orthonormées ; coordonnées comme produits scalaires & 16\\
Normes d'une somme et d'une différence ; Al-Kashi ; calculs de longueurs et d'angles & 16--17\\
Transformer $\overrightarrow{MA}\cdot\overrightarrow{MB}$ ; cercle de diamètre $[AB]$ et cas dégénérés & 17\\
Vecteur normal, équation $ax+by+c=0$, droites perpendiculaires ; projeté orthogonal d'un point & 18 ; S3\\
Équation d'un cercle, lecture centre/rayon ; problèmes de configurations en repère & 18 ; S3\\
Approfondissements : sinus, médianes et hauteurs concourantes, parabole foyer/directrice, intersections & 17--18\\\bottomrule
\end{longtable}
\section{Probabilités et statistiques}
\begin{longtable}{>{\raggedright\arraybackslash}p{.70\linewidth}>{\raggedright\arraybackslash}p{.20\linewidth}}
\toprule Contenus et capacités & Localisation\\\midrule
Indépendance de deux événements ; complémentaire ; conditionnement et différence avec intersection & 19 ; guide\\
Partition et formule des probabilités totales ; tableaux croisés et arbres pondérés & 19 ; guide ; S4\\
Deux expériences indépendantes, épreuves répétées de Bernoulli, jusqu'à quatre épreuves dans le socle & 20 ; S5\\
Univers fini, variable aléatoire, événements $\{X=a\}$ et $\{X\le a\}$, loi de probabilité & 21\\
Espérance, variance, écart type ; linéarité de l'espérance et König--Huygens ; jeux et mises & 21--22 ; S5\\
Simulation d'une variable, moyenne d'un échantillon ; répétition de $N$ échantillons de taille $n$ ; proportion dans $2\sigma/\sqrt n$ & 23--24 ; Python\\
Monte-Carlo : aire sous une parabole et approximation de $\pi$ ; fréquences de lettres & 23--24 ; Python\\
Approfondissements : loi binomiale générale, marche aléatoire, covariance, erreur quadratique minimale & 20, 22, 28\\
Réactivation statistique : moyenne pondérée, médiane, quartiles, étendue, écart interquartile, variance ; représentations & guide ; seconde 19--20\\\bottomrule
\end{longtable}
\section{Trois dimensions transversales}
\begin{longtable}{>{\raggedright\arraybackslash}p{.70\linewidth}>{\raggedright\arraybackslash}p{.20\linewidth}}
\toprule Contenus et capacités & Localisation\\\midrule
Ensembles : appartenance, inclusion, vide, réunion, intersection, complémentaire, couples, produit cartésien, cardinal & 1\\
Logique : et/ou, implication, équivalence, réciproque, contraposée, nécessaire/suffisant, quantification en français et négation & 1 et preuves\\
Variables, inconnues, paramètres, identités et équations ; raisonnements par cas, absurde, contraposée & 1--3, 7--10, 25\\
Algorithmique : types, affectations, conditions, boucles, fonctions, programmation modulaire & 4, 24 ; Python\\
Listes en extension/compréhension, ajout, suppression, indices, parcours et agrégation & 24 ; Python\\
Automatismes : calcul numérique et littéral, équations, signes, taux directs/réciproques/successifs, données, fonctions et droites & guide, 80 questions\\
Chercher, modéliser, représenter, raisonner, calculer, communiquer ; outils numériques et oral & tous les chapitres ; S1 à S8\\\bottomrule
\end{longtable}

\chapter{Démonstrations et fondements}
\section{Démonstrations identifiées dans le programme}
La liste ci-dessous permet de retrouver les preuves exemplaires explicitement
indiquées par le programme. Elle ne limite pas les démonstrations du cours.
\begin{longtable}{>{\raggedright\arraybackslash}p{.72\linewidth}>{\raggedright\arraybackslash}p{.18\linewidth}}
\toprule Résultat & Chapitre\\\midrule
Terme général d'une suite arithmétique et géométrique & 5\\
Sommes des premiers entiers et des puissances d'un réel & 5\\
Formules des racines d'un polynôme de degré deux & 3\\
Équation de la tangente au point d'abscisse $a$ & 8\\
Non-dérivabilité finie de la racine carrée en zéro & 8\\
Dérivées des fonctions carré et inverse & 8--9\\
Dérivée d'un produit & 9\\
Valeurs de sinus et cosinus en $\pi/4$ et $\pi/3$ & 14\\
Théorème d'Al-Kashi & 17\\
Lieu $\overrightarrow{MA}\cdot\overrightarrow{MB}=0$ & 17\\\bottomrule
\end{longtable}
Le cours démontre aussi : quotient et puissances, Fermat, Rolle et accroissements
finis, règles de l'exponentielle, Bayes, probabilités totales, linéarité et variance,
Cauchy--Schwarz, projections, variance totale, Tchebychev, contractions et erreurs
de Newton et d'Euler.

\section{Dépendances et admissions clairement identifiées}
\begin{longtable}{>{\raggedright\arraybackslash}p{.36\linewidth}>{\raggedright\arraybackslash}p{.54\linewidth}}
\toprule Outil & Fondement ou preuve disponible\\\midrule
Calcul réel, ordre, propriétés usuelles des entiers & Fondements algébriques admis ; acquis de collège et de seconde.\\
Récurrence sur les entiers & Principe de base ; premières preuves au chapitre 1.\\
Pythagore, Thalès, triangles semblables et géométrie euclidienne & Acquis démontrés dans les volumes précédents ; utilisés explicitement ici.\\
Longueur d'arc, aire du disque/secteur, comparaison polygonale & Fondements géométriques utilisés au chapitre 14 et dans la limite trigonométrique ; une construction complète des longueurs courbes n'est pas entreprise.\\
Complétude de $\R$ & Borne supérieure admise au chapitre 25 ; convergence monotone déduite au chapitre 7.\\
Bornes atteintes pour une fonction continue & Démontrées au chapitre 25, avant Rolle et les accroissements finis du chapitre 10.\\
Existence de l'exponentielle & Admise dans le socle, construite par séries et intégration au chapitre 26 ; unicité prouvée au chapitre 12.\\
Intégrale d'une fonction continue & Construite au chapitre 25 par sommes inférieures/supérieures et continuité uniforme.\\
Modèle uniforme continu dans un carré & Probabilités définies par les aires pour Monte-Carlo ; cadre explicitement plus large que l'univers fini du socle.\\
Fonctionnement des flottants et du générateur Python & Modèles informatiques utilisés avec leurs limites ; les preuves concernent l'arithmétique mathématique et l'indépendance hypothétique.\\\bottomrule
\end{longtable}
Les chapitres 15 et 16 se lisent avec un renvoi en avant : les formules
trigonométriques utilisent le produit scalaire, mais la preuve du produit scalaire
ne suppose pas ces formules. Les chapitres 25 et 26 ferment les dépendances
annoncées plus tôt ; les preuves ne supposent pas leurs propres conclusions.

\chapter{Réactiver les acquis et lire des données}
\section{Pourcentages et changements de référence}
Un taux $t$ transforme $V_i$ en $V_f=(1+t)V_i$.
Pour deux évolutions, le coefficient global est $(1+t_1)(1+t_2)$ ;
les taux ne s'additionnent pas en général.
Le taux réciproque vaut $1/(1+t)-1=-t/(1+t)$ lorsque $1+t>0$.
Ces formules viennent de la multiplication puis de l'inversion du coefficient.
Par exemple, une hausse de $25\%$ s'annule par une baisse de $20\%$.
Une proportion de proportion se multiplie : $30\%$ de $40\%$ représente
$0{,}3\times0{,}4=0{,}12$, soit $12\%$ de l'ensemble initial.

\section{Indicateurs et conventions}
Pour une série de valeurs $x_i$ d'effectifs $n_i$, $N=\sum n_i$,
la moyenne est $\overline x=\sum n_ix_i/N$.
La variance descriptive est $\sum n_i(x_i-\overline x)^2/N$.
Les démonstrations de König--Huygens et du changement d'unité sont au chapitre 22.
Une médiane partage la série ordonnée en deux parties d'au moins la moitié
des observations ; pour un effectif pair, on choisit ici la moyenne des deux
valeurs centrales. Les quartiles suivent ici la convention française des rangs
$\lceil N/4\rceil$ et $\lceil3N/4\rceil$.
L'étendue est maximum moins minimum ; l'écart interquartile est $Q_3-Q_1$.

Pour $1,2,2,3,4,4,4,8$, la moyenne vaut $28/8=3{,}5$,
la médiane $(3+4)/2=3{,}5$, $Q_1=2$, $Q_3=4$, l'étendue $7$
et l'écart interquartile $2$.
La somme des carrés est $130$ ; la variance vaut $130/8-(3{,}5)^2=4$,
donc l'écart type vaut $2$.
Une boîte à moustaches peut représenter ces cinq repères ; les moustaches
sont ici le minimum et le maximum, non une convention d'exclusion des valeurs extrêmes.
\begin{center}
\begin{tikzpicture}[xscale=.8]
\draw[->](0,0)--(9,0);
\foreach \x in {1,2,3,4,5,6,7,8}{\draw(\x,.08)--(\x,-.08)node[below]{\x};}
\draw(1,1)--(2,1) (4,1)--(8,1);
\draw(1,.7)--(1,1.3) (8,.7)--(8,1.3);
\draw[fill=accent!12](2,.6)rectangle(4,1.4);
\draw[thick](3.5,.6)--(3.5,1.4);
\end{tikzpicture}
\end{center}

\section{Histogramme, bâtons et secteurs}
Pour une variable discrète, un diagramme en bâtons place chaque valeur en abscisse
et son effectif ou sa fréquence en hauteur. Pour des classes, un histogramme
représente les effectifs par les \emph{aires} des rectangles.
Avec $10$ observations dans $[0,10[$ et $30$ dans $[10,30]$,
les hauteurs en effectifs par unité sont $10/10=1$ et $30/20=1{,}5$.
Les classes ont des largeurs différentes : dessiner des hauteurs $10$ et $30$
fausserait le rapport des aires.
\begin{center}
\begin{tikzpicture}
\begin{axis}[width=.65\linewidth,height=4.5cm,axis lines=left,
xmin=0,xmax=32,ymin=0,ymax=1.8,xtick={0,10,30},ytick={1,1.5},
xlabel={Valeur},ylabel={Densité d'effectif}]
\addplot[fill=accent!15,draw=accent]coordinates{(0,0)(0,1)(10,1)(10,0)}\closedcycle;
\addplot[fill=accent!35,draw=accent]coordinates{(10,0)(10,1.5)(30,1.5)(30,0)}\closedcycle;
\end{axis}
\end{tikzpicture}
\end{center}
La moyenne exacte d'une série groupée en classes n'est généralement pas
déterminée ; remplacer chaque classe par son centre fournit une approximation.
Dans un diagramme circulaire, une proportion $p$ correspond à un angle
$360p$ degrés, par proportionnalité à un tour entier.

\section{Lire une courbe, une droite et une tangente}
Pour lire $f(a)$, partir de l'abscisse $a$, rejoindre la courbe puis lire
l'ordonnée. Pour résoudre $f(x)=k$, chercher les intersections avec $y=k$ ;
pour $f(x)\le k$, repérer les abscisses où la courbe est sous cette droite.
Le signe de $f$ concerne la position par rapport à l'axe horizontal ;
les variations concernent l'évolution des ordonnées quand l'abscisse augmente.
Une fonction peut être positive et décroissante.

Une droite passant par $A(x_A,y_A)$ et $B(x_B,y_B)$ avec $x_A\ne x_B$
a pour pente $(y_B-y_A)/(x_B-x_A)$, obtenue en divisant les variations.
Son équation est $y-y_A=m(x-x_A)$ ; si $x_A=x_B$, elle est verticale
d'équation $x=x_A$. La pente d'une tangente se lit avec deux points de
\emph{la tangente}, qui ne sont pas forcément deux points de la courbe.
La figure du chapitre 8 donne un exemple exact : pour $f(x)=x^2$,
la tangente en $1$ est $y=2x-1$.

\section{Arbre pondéré et tableau}
Avec une partition $A,\overline A$, les premières branches portent
$P(A)$ et $P(\overline A)$ ; les branches suivantes portent des
probabilités conditionnelles.
\begin{center}
\begin{tikzpicture}[grow=right,level distance=3.2cm,
level 1/.style={sibling distance=3.1cm},
level 2/.style={sibling distance=1.5cm},
edge from parent/.style={draw,-Latex}]
\node {}
child {node {$\overline A$}
 child {node {$\overline B$} edge from parent node[below,sloped] {$1-q$}}
 child {node {$B$} edge from parent node[above,sloped] {$q$}}
 edge from parent node[below,sloped] {$1-p$}}
child {node {$A$}
 child {node {$\overline B$} edge from parent node[below,sloped] {$1-r$}}
 child {node {$B$} edge from parent node[above,sloped] {$r$}}
 edge from parent node[above,sloped] {$p$}};
\end{tikzpicture}
\end{center}
Ici $r=P_A(B)$ et $q=P_{\overline A}(B)$.
Les chemins vers $B$ ont probabilités $pr$ et $(1-p)q$,
d'où $P(B)=pr+(1-p)q$. Dans un tableau, ces chemins sont les cases
$A\cap B$ et $\overline A\cap B$ ; conditionner revient à diviser par
le total de la ligne ou de la colonne correspondante.

\chapter{Quatre-vingts automatismes}
Ces questions courtes entretiennent les acquis de seconde et les nouveaux outils.
Les questions 71 à 80 sont des rappels de preuve ou des compléments signalés.
Travailler par groupes de dix, sans consulter immédiatement les réponses.
\section{Questions 1 à 10}
\begin{enumerate}[start=1]
\item Calculer $3/4-5/6$.
\item Simplifier $\sqrt{72}$.
\item Calculer $2^{-3}$.
\item Écrire $0{,}00042$ en notation scientifique.
\item Développer $(x-4)^2$.
\item Factoriser $9x^2-25$.
\item Simplifier $(x^2-4)/(x-2)$ et donner le domaine.
\item Résoudre $3x-7=5$.
\item Résoudre $-2x+1>7$.
\item Résoudre $|x-2|\le3$.
\end{enumerate}
\section{Questions 11 à 20}
\begin{enumerate}[start=11]
\item Calculer $15\%$ de $240$.
\item Une valeur de $80$ augmente de $25\%$. Valeur finale ?
\item Après une baisse de $20\%$, la valeur est $64$. Valeur initiale ?
\item Taux global d'une hausse de $10\%$ suivie d'une baisse de $10\%$ ?
\item Taux réciproque d'une hausse de $25\%$ ?
\item $40\%$ d'une population, dont $30\%$ ont une propriété : proportion totale ?
\item De $50$ à $65$, quel taux ?
\item Deux hausses de $20\%$ donnent quel coefficient ?
\item Un coefficient global $0{,}81$ après deux baisses identiques : taux par étape ?
\item Une proportion passe de $20\%$ à $25\%$. Hausse en points et hausse relative ?
\end{enumerate}
\section{Questions 21 à 30}
\begin{enumerate}[start=21]
\item Racines de $x^2-5x+6$ ?
\item Signe de $(x+1)(x-4)$ ?
\item Forme canonique de $x^2+6x+5$ ?
\item Discriminant de $2x^2+4x+3$ ?
\item Somme et produit des racines de $3x^2-6x-2$, qui possède deux racines ?
\item Résoudre $x^2<9$.
\item Minimum de $2(x-3)^2+1$ ?
\item Domaine de $\sqrt{2x-4}$ ?
\item Résoudre $1/(x-1)=2$.
\item Le polynôme nul a-t-il au plus deux racines ?
\end{enumerate}
\section{Questions 31 à 40}
\begin{enumerate}[start=31]
\item Suite arithmétique $u_0=4$, raison $3$ : $u_{10}$ ?
\item Suite géométrique $u_0=5$, raison $2$ : $u_4$ ?
\item Calculer $1+\cdots+20$.
\item Calculer $1+2+4+\cdots+2^6$.
\item Nombre de termes de $u_3$ à $u_{12}$ inclus ?
\item Équilibre de $u_{n+1}=0{,}5u_n+6$ ?
\item Limite de $3+2(0{,}9)^n$ ?
\item La suite $(-1)^n$ est-elle convergente ?
\item Si $u_{n+1}-u_n>0$, quelle monotonie ?
\item Que renvoie \texttt{list(range(2, 5))} ?
\end{enumerate}
\section{Questions 41 à 50}
\begin{enumerate}[start=41]
\item Dérivée de $x^3-4x$ ?
\item Dérivée de $1/x$ ?
\item Dérivée de $\sqrt x$ ?
\item Dérivée de $(2x+1)^3$ ?
\item Tangente à $x^2$ en $2$ ?
\item $f'(x)<0$ sur un intervalle : quelle variation ?
\item $f'(a)=0$ garantit-il un extremum ?
\item Dérivée de $e^{3x}$ ?
\item Résoudre $e^{2x}<e^6$.
\item Simplifier $e^{x+1}e^{-x}$.
\end{enumerate}
\section{Questions 51 à 60}
\begin{enumerate}[start=51]
\item Convertir $150^\circ$ en radians.
\item Valeur de $\sin(\pi/3)$ ?
\item Valeur de $\cos(3\pi/4)$ ?
\item Calculer $(2,-1)\cdot(3,4)$.
\item Norme de $(3,-4)$ ?
\item Un normal à $2x-3y+5=0$ ?
\item Pente de la droite passant par $(1,2)$ et $(3,8)$ ?
\item Équation de cette droite ?
\item Centre et rayon de $(x+2)^2+(y-1)^2=9$ ?
\item Si $\|u\|=2$, $\|v\|=3$, $u\perp v$, calculer $\|u+v\|$.
\end{enumerate}
\section{Questions 61 à 70}
\begin{enumerate}[start=61]
\item Moyenne de $2,3,3,8$ ?
\item Médiane et étendue de $1,2,4,7,9$ ?
\item Une classe de largeur $5$ contient $20$ données. Hauteur d'un histogramme d'effectifs ?
\item Angle d'un secteur représentant $15\%$ ?
\item $P(A)=0{,}4$ et $P_A(B)=0{,}3$ : $P(A\cap B)$ ?
\item $P(A\cap B)=0{,}12$ et $P(B)=0{,}6$ : $P_B(A)$ ?
\item Deux événements indépendants de probabilités $0{,}2$ et $0{,}7$ : intersection ?
\item Espérance d'une Bernoulli de paramètre $p$ ?
\item Variance de cette Bernoulli ?
\item Écart type d'une moyenne de $100$ variables indépendantes de même écart type $5$ ?
\end{enumerate}
\section{Questions 71 à 80}
\begin{enumerate}[start=71]
\item Nier « pour tout $x$, $f(x)>0$ ».
\item Contraposée de « $A$ implique $B$ » ?
\item Une propriété vraie sur mille exemples est-elle démontrée sur $\N$ ?
\item Une indépendance est-elle nécessaire à $\mathbb E(X+Y)=\mathbb E(X)+\mathbb E(Y)$ ?
\item Complément : dérivée de $\ln x$ ?
\item Complément : dérivée de $\sin x$ en radians ?
\item Complément : meilleure constante pour approcher $X$ en erreur quadratique ?
\item Complément : projection de $(2,3)$ sur l'axe des abscisses ?
\item Complément : une contraction de rapport $q<1$ peut-elle avoir deux points fixes ?
\item Une fréquence simulée est-elle une probabilité exacte ?
\end{enumerate}
\chapter{Réponses et contrôles}
\section{Réponses 1 à 10}
\begin{enumerate}[start=1]
\item $-1/12$, après mise au dénominateur $12$.
\item $6\sqrt2$, car $72=36\times2$.
\item $1/8$.
\item $4{,}2\times10^{-4}$.
\item $x^2-8x+16$.
\item $(3x-5)(3x+5)$.
\item $x+2$ pour $x\ne2$.
\item $x=4$.
\item $x<-3$ : diviser par un négatif inverse le sens.
\item $-1\le x\le5$.
\end{enumerate}
\section{Réponses 11 à 20}
\begin{enumerate}[start=11]
\item $36$.
\item $80\times1{,}25=100$.
\item $64/0{,}8=80$.
\item $1{,}1\times0{,}9-1=-0{,}01$, baisse de $1\%$.
\item $1/1{,}25-1=-20\%$.
\item $0{,}4\times0{,}3=12\%$.
\item $(65-50)/50=30\%$.
\item $1{,}2^2=1{,}44$.
\item $\sqrt{0{,}81}=0{,}9$, soit une baisse de $10\%$.
\item $5$ points ; hausse relative $(25-20)/20=25\%$.
\end{enumerate}
\section{Réponses 21 à 30}
\begin{enumerate}[start=21]
\item $2$ et $3$, par $(x-2)(x-3)$.
\item Positif hors de $[-1,4]$, négatif sur $]-1,4[$, nul aux bornes.
\item $(x+3)^2-4$.
\item $16-24=-8$.
\item Somme $2$, produit $-2/3$.
\item $-3<x<3$.
\item $1$, atteint en $x=3$.
\item $[2,+\infty[$.
\item $x=3/2$, avec $x\ne1$.
\item Non : tout réel est racine ; le théorème exige un polynôme non nul de degré au plus deux.
\end{enumerate}
\section{Réponses 31 à 40}
\begin{enumerate}[start=31]
\item $4+10\times3=34$.
\item $5\times2^4=80$.
\item $20\times21/2=210$.
\item $(2^7-1)/(2-1)=127$.
\item $12-3+1=10$.
\item $12$, solution de $\ell=0{,}5\ell+6$.
\item $3$.
\item Non, ses sous-suites paires et impaires ont des limites différentes.
\item Suite strictement croissante.
\item \texttt{[2, 3, 4]}.
\end{enumerate}
\section{Réponses 41 à 50}
\begin{enumerate}[start=41]
\item $3x^2-4$.
\item $-1/x^2$ pour $x\ne0$.
\item $1/(2\sqrt x)$ pour $x>0$.
\item $6(2x+1)^2$.
\item $y=4+4(x-2)=4x-4$.
\item $f$ est strictement décroissante sur cet intervalle.
\item Non : $x^3$ en $0$ est un contre-exemple.
\item $3e^{3x}$.
\item $x<3$ par stricte croissance.
\item $e$.
\end{enumerate}
\section{Réponses 51 à 60}
\begin{enumerate}[start=51]
\item $5\pi/6$.
\item $\sqrt3/2$.
\item $-\sqrt2/2$.
\item $6-4=2$.
\item $5$.
\item $(2,-3)$.
\item $(8-2)/(3-1)=3$.
\item $y=3x-1$.
\item Centre $(-2,1)$, rayon $3$.
\item $\sqrt{4+9}=\sqrt{13}$.
\end{enumerate}
\section{Réponses 61 à 70}
\begin{enumerate}[start=61]
\item $4$.
\item Médiane $4$, étendue $8$.
\item $20/5=4$ observations par unité.
\item $0{,}15\times360^\circ=54^\circ$.
\item $0{,}12$.
\item $0{,}2$.
\item $0{,}14$.
\item $p$.
\item $p(1-p)$.
\item $5/\sqrt{100}=0{,}5$.
\end{enumerate}
\section{Réponses 71 à 80}
\begin{enumerate}[start=71]
\item « Il existe $x$ tel que $f(x)\le0$ », sur le même domaine.
\item « Non $B$ implique non $A$ ».
\item Non : il faut un argument couvrant tous les cas.
\item Non, la linéarité vaut pour toutes les variables finies définies ensemble.
\item $1/x$ pour $x>0$.
\item $\cos x$.
\item $\mathbb E(X)$, par la décomposition variance plus carré.
\item $(2,0)$ ; le résidu $(0,3)$ est orthogonal à l'axe.
\item Non : leur distance serait au plus $q$ fois elle-même, donc nulle.
\item Non : elle fluctue ; taille, modèle et méthode déterminent comment l'interpréter.
\end{enumerate}

\chapter{Préparer les études supérieures}
\section{Des outils et des usages précis}
\begin{longtable}{>{\raggedright\arraybackslash}p{.25\linewidth}>{\raggedright\arraybackslash}p{.35\linewidth}>{\raggedright\arraybackslash}p{.30\linewidth}}
\toprule Outil actuel & Premier réemploi démontré ici & Développement ultérieur\\\midrule
Identités, factorisation & Racines, signes, sommes télescopiques & Polynômes, structures algébriques\\
Suites géométriques & Modèles affines, stabilité, contractions & Séries, systèmes dynamiques\\
Limites et complétude & Existence d'un point fixe, bornes atteintes & Topologie, espaces fonctionnels\\
Dérivée et tangente & Optimisation, Newton, approximation affine & Calcul différentiel, gradients\\
Exponentielle & Modèles continus, Euler, logarithme & Équations différentielles, matrices\\
Trigonométrie & Rotations, oscillations, signal à quatre points & Complexes, Fourier, analyse harmonique\\
Produit scalaire & Projection, moindres carrés & Algèbre linéaire, espaces de Hilbert\\
Espérance et variance & Concentration, prédiction conditionnelle & Mesure, statistiques, probabilités avancées\\
Listes et invariants & Algorithmes vérifiés, erreurs contrôlées & Calcul scientifique et optimisation\\\bottomrule
\end{longtable}

\section{Maîtrise et relecture}
Avant de passer à un nouvel objet, savoir expliquer l'ancien sans réciter :
définition, exemple, contre-exemple, théorème avec hypothèses, preuve et usage.
Les compléments demandent un effort nouveau ; les aborder progressivement,
puis les reprendre après les futurs volumes de terminale et de licence.
Le chapitre 30 donne une progression indicative jusqu'au master, dont la
spécialisation devra être choisie ultérieurement.

\section{Reproduire les calculs et les documents}
Depuis la racine du dépôt :
\begin{lstlisting}
python3 compiler.py --niveau premiere
python3 compiler.py --niveau premiere --chapitre 16
python3 premiere/python/verifier.py
python3 premiere/python/outils.py
\end{lstlisting}
Le style \texttt{mathscommun.sty} est commun aux niveaux. Le module Python précise
les conventions numériques et les graines des simulations.
\end{document}
